% GENERATED FILE. Edit canonical lessons, then run npm run generate:books.
% canonical-source-hash: fnv1a64:e03b1e882ebaa629
% canonical-lessons: BN-C149-satar-kata, BN-C149-pichhle-pora, BN-C149-bishram-neoya, BN-C149-taratari-kora, BN-C149-deri-kora

\chapter{To swim, To slip, To rest, To hurry, To delay}
\label{ch:bn-to-swim-to-slip-to-rest-to-hurry-to-delay}
\hlchaptermodality{149}

\begin{chapteropening}
\textbf{By the end of this chapter:} \emph{I can say to swim, to slip, to rest, to hurry and to delay.}

Complete the last lesson of chapter 149: I can say to swim, to slip, to rest, to hurry and to delay.
\end{chapteropening}

\section[\texorpdfstring{\bn{সাঁতার} \bn{কাটা} (sā̃tār kāṭā)}{সাঁতার কাটা (sā̃tār kāṭā)}]{\bn{সাঁতার} \bn{কাটা} (sā̃tār kāṭā) — to swim}

\label{lesson:BN-C149-satar-kata}



\begin{quote}
Before the new one: say the Bengali for to sweep, then the Bengali for to boil.
\end{quote}

\subsection*{You'll want to know: \bn{সাঁতার} \bn{কাটা}}
\textbf{\bn{সাঁতার} \bn{কাটা}} — \emph{sā̃tār kāṭā} — \textquotedblleft{}to swim\textquotedblright{}.

\begin{grammarlens}[title={What you've built}]
One more word for this chapter.
\end{grammarlens}

\subsection*{Your turn}
\begin{itemize}
\raggedright
  \item \emph{Say it:} \emph{sā̃tār kāṭā}
  \item \emph{Say it:} \emph{sā̃tār kāṭā}, once more
  \item \emph{Say it:} say \emph{sā̃tār kāṭā}
  \item \emph{Recall:} say the Bengali for to sweep, then the Bengali for to boil, then say \emph{sā̃tār kāṭā} again
\end{itemize}

\begin{tcolorbox}[breakable,colback=teal!4,colframe=teal!35!black,title={Before you move on}]
What does \bn{সাঁতার} \bn{কাটা} mean? (\textquotedblleft{}To swim\textquotedblright{}.) Say it once more.
\end{tcolorbox}
\section[\texorpdfstring{\bn{পিছলে} \bn{পড়া} (pichhle pôṛā)}{পিছলে পড়া (pichhle pôṛā)}]{\bn{পিছলে} \bn{পড়া} (pichhle pôṛā) — to slip}

\label{lesson:BN-C149-pichhle-pora}



\begin{quote}
Before the new one: say the Bengali for to boil, then the Bengali for to swim.
\end{quote}

\subsection*{You'll want to know: \bn{পিছলে} \bn{পড়া}}
\textbf{\bn{পিছলে} \bn{পড়া}} — \emph{pichhle pôṛā} — \textquotedblleft{}to slip\textquotedblright{}.

\begin{grammarlens}[title={What you've built}]
One more word for this chapter.
\end{grammarlens}

\subsection*{Your turn}
\begin{itemize}
\raggedright
  \item \emph{Say it:} \emph{pichhle pôṛā}
  \item \emph{Say it:} \emph{pichhle pôṛā}, once more
  \item \emph{Say it:} say \emph{pichhle pôṛā}
  \item \emph{Recall:} say the Bengali for to boil, then the Bengali for to swim, then say \emph{pichhle pôṛā} again
\end{itemize}

\begin{tcolorbox}[breakable,colback=teal!4,colframe=teal!35!black,title={Before you move on}]
What does \bn{পিছলে} \bn{পড়া} mean? (\textquotedblleft{}To slip\textquotedblright{}.) Say it once more.
\end{tcolorbox}
\section[\texorpdfstring{\bn{বিশ্রাম} \bn{নেওয়া} (bishrām neoyā)}{বিশ্রাম নেওয়া (bishrām neoyā)}]{\bn{বিশ্রাম} \bn{নেওয়া} (bishrām neoyā) — to rest}

\label{lesson:BN-C149-bishram-neoya}



\begin{quote}
Before the new one: say the Bengali for to swim, then the Bengali for to slip.
\end{quote}

\subsection*{You'll want to know: \bn{বিশ্রাম} \bn{নেওয়া}}
\textbf{\bn{বিশ্রাম} \bn{নেওয়া}} — \emph{bishrām neoyā} — \textquotedblleft{}to rest\textquotedblright{}.

\begin{grammarlens}[title={What you've built}]
One more word for this chapter.
\end{grammarlens}

\subsection*{Your turn}
\begin{itemize}
\raggedright
  \item \emph{Say it:} \emph{bishrām neoyā}
  \item \emph{Say it:} \emph{bishrām neoyā}, once more
  \item \emph{Say it:} say \emph{bishrām neoyā}
  \item \emph{Recall:} say the Bengali for to swim, then the Bengali for to slip, then say \emph{bishrām neoyā} again
\end{itemize}

\begin{tcolorbox}[breakable,colback=teal!4,colframe=teal!35!black,title={Before you move on}]
What does \bn{বিশ্রাম} \bn{নেওয়া} mean? (\textquotedblleft{}To rest\textquotedblright{}.) Say it once more.
\end{tcolorbox}
\section[\texorpdfstring{\bn{তাড়াতাড়ি} \bn{করা} (tāṛātāṛi kôrā)}{তাড়াতাড়ি করা (tāṛātāṛi kôrā)}]{\bn{তাড়াতাড়ি} \bn{করা} (tāṛātāṛi kôrā) — to hurry}

\label{lesson:BN-C149-taratari-kora}



\begin{quote}
Before the new one: say the Bengali for to slip, then the Bengali for to rest.
\end{quote}

\subsection*{You'll want to know: \bn{তাড়াতাড়ি} \bn{করা}}
\textbf{\bn{তাড়াতাড়ি} \bn{করা}} — \emph{tāṛātāṛi kôrā} — \textquotedblleft{}to hurry\textquotedblright{}.

\begin{grammarlens}[title={What you've built}]
One more word for this chapter.
\end{grammarlens}

\subsection*{Your turn}
\begin{itemize}
\raggedright
  \item \emph{Say it:} \emph{tāṛātāṛi kôrā}
  \item \emph{Say it:} \emph{tāṛātāṛi kôrā}, once more
  \item \emph{Say it:} say \emph{tāṛātāṛi kôrā}
  \item \emph{Recall:} say the Bengali for to slip, then the Bengali for to rest, then say \emph{tāṛātāṛi kôrā} again
\end{itemize}

\begin{tcolorbox}[breakable,colback=teal!4,colframe=teal!35!black,title={Before you move on}]
What does \bn{তাড়াতাড়ি} \bn{করা} mean? (\textquotedblleft{}To hurry\textquotedblright{}.) Say it once more.
\end{tcolorbox}
\section[\texorpdfstring{\bn{দেরি} \bn{করা} (deri kôrā)}{দেরি করা (deri kôrā)}]{\bn{দেরি} \bn{করা} (deri kôrā) — to delay, to be late}

\label{lesson:BN-C149-deri-kora}



\begin{quote}
Before the new one: say the Bengali for to rest, then the Bengali for to hurry.
\end{quote}

\subsection*{You'll want to know: \bn{দেরি} \bn{করা}}
\textbf{\bn{দেরি} \bn{করা}} — \emph{deri kôrā} — \textquotedblleft{}to delay, to be late\textquotedblright{}.

\begin{grammarlens}[title={What you've built}]
One more word for this chapter.
\end{grammarlens}

\subsection*{Your turn}
\begin{itemize}
\raggedright
  \item \emph{Say it:} \emph{deri kôrā}
  \item \emph{Say it:} \emph{deri kôrā}, once more
  \item \emph{Say it:} say \emph{deri kôrā}
  \item \emph{Recall:} say the Bengali for to rest, then the Bengali for to hurry, then say \emph{deri kôrā} again
\end{itemize}

\begin{tcolorbox}[breakable,colback=teal!4,colframe=teal!35!black,title={Before you move on}]
What does \bn{দেরি} \bn{করা} mean? (\textquotedblleft{}To delay, to be late\textquotedblright{}.) Say it once more.
\end{tcolorbox}
